% ===== tr-04 块2：第 4 章 pp.14--26；本片 pp.14--18（原文页 14--18）=====
% [原文 p.14]
我们在此不给出这是如何完成的具体细节，但它并不特别困难，只要把范围限制在拓扑（而非 $C^\infty$）流形的范畴内。不过，看一看我们需要添入的奇异纤维是很有意思的。

若 $b\in\Gamma$ 是 $\Gamma$ 中一个非顶点，则 $f^{-1}(b)$ 同胚于 $I_1\times S^1$，其中 $I_1$ 表示 Kodaira 类型 $I_1$ 的椭圆曲线，即尖点环面（pinched torus）。
% [Tutor Note: 此处的 I_1 是 Kodaira 纤维类型记号（尖点椭圆曲线，即把环面的一条 S^1 捏成一点），不是实区间 [0,1]，也不是某个指标 i；I_1 × S^1 与三维尖点环面同胚。]

若 $v$ 是 $\Gamma$ 的一个顶点、单值为第一种类型，则 $f^{-1}(v)=S^1\times S^1\times S^1/{\sim}$，其中 $(a,b,c)\sim(a',b',c')$ 当且仅当 $(a,b,c)=(a',b',c')$ 或 $a=a'=1$；这里 $S^1$ 等同于 $\C$ 中的单位圆。这是尖点环面的三维类比，且 $\chi(f^{-1}(v))=+1$。我们称它为一个\emph{正纤维}。

若 $v$ 是 $\Gamma$ 的一个顶点、单值为第二种类型，则 $f^{-1}(v)$ 可以描述为 $S^1\times S^1\times S^1/{\sim}$，其中 $(a,b,c)\sim(a',b',c')$ 当且仅当 $(a,b,c)=(a',b',c')$，或 $a=a'=1$ 且 $b=b'$，或 $a=a'$ 且 $b=b'=1$。该纤维的奇异轨迹是一个八字形，且 $\chi(f^{-1}(v))=-1$。我们称它为一个\emph{负纤维}。
% [Tutor Note: 正/负纤维的差别只在等价关系：正型只有 a=a'=1 一个附加条件（坍缩一个 S^1），负型有 a=a'=1,b=b' 与 a=a',b=b'=1 两个附加条件（坍缩两条 S^1）；Euler 示性数分别为 +1 与 −1，抄写时易把两个等价关系合并。]

于是我们看到 SYZ 对偶的一个非常具体的局部推论：在紧化 $X(B)$ 与 $\check X(B)$ 中，正纤维与负纤维互换。当然，这导致如下观察：对 Calabi--Yau 三维流形，Euler 示性数在镜像对称下变号。

\begin{example}{4.3.4}
接续例 3.2，[20] 中证明了 $\check X(B)$ 同胚于五次超曲面，而 $X(B)$ 同胚于镜像五次超曲面。除去一篇因其他更紧迫的项目而长期延迟问世的论文 [30]，[22] 的结果蕴含 SYZ 猜想在拓扑层面对环面簇中的所有完全交成立。
\end{example}

W.-D. Ruan 在 [70] 中用辛流论证给出了环面簇中超曲面的拉格朗日环面纤维化的一个描述，他的构造应当与辛流形 $\check X(B_0)$ 的辛紧化一致。在三维情形，这样的辛紧化已由 Ricardo Castaño-Bernard 与 Diego Matessi [8] 构造出来。若把这一紧化应用于此处描述的带奇点仿射流形，所得的辛流形应当与相应的环面超曲面辛同胚，但这尚未被证明。

\section{热带几何}
回忆镜像对称理应允许我们数曲线，让我们在直观层面讨论一下目前的图景如何让我们洞悉这个问题。设 $B$ 是一个热带仿射流形。则如我们所见，$X(B)$ 带有半平坦复结构，而 $X(B)$ 的某些复子流形可以很容易地如下描述。设 $L\subseteq B$ 是一个
% [原文 p.15]
具有有理斜率的线性子空间，即对任意 $b\in L$，$L$ 在 $b$ 处的切空间 $\mathcal{T}_{L,b}$ 可以写成 $M\otimes_{\Z}\R$，其中 $M\subseteq\Lambda_b$ 是某个子格。于是我们得到子流形
\[
X(L):=\mathcal{T}_L/(\mathcal{T}_L\cap\Lambda)\subseteq X(B).
\]
容易验证这是一个复子流形。例如，若 $B=\R^n$，从而 $X(B)$ 只是一个代数环面 $(\C^*)^n$，其坐标为 $q_1,\dots,q_n$，而 $L\subseteq B$ 是由方程
\[
\sum_j c_{ij}y_j=d_i,\qquad 1\le i\le p
\]
定义的余维数 $p$ 的仿射线性子空间，其中 $c_{ij}\in\Z$、$d_i\in\R$，则 $X(B)$ 中相应的子流形是由方程
\[
\prod_j q_j^{c_{ij}}=e^{-2\pi d_j},\qquad 1\le i\le p
\]
给出的子环面。
% [Tutor Note: 此式右端原文印作 e^{-2πd_j}，但其左侧与约束 1≤i≤p 都是第 i 个方程，按 ∏_j q_j^{c_{ij}} = e^{-2π Σ_j c_{ij}y_j} 应为 e^{-2πd_i}。已 900/1200 dpi 裁图核实原文字母确为 d_j，正文照录，判为原文笔误。]

当然，环面的子环面并不特别有意思。我们怎样构造更复杂的子流形呢？让我们聚焦于曲线，即把 $B$ 的线性子流形取为一维的。于是若取 $L$ 为线段、射线或直线，$X(L)$ 就是一个圆柱面，在各种情形下带或不带边界。然后我们可以试着把这样的圆柱面粘合起来，得到更复杂的曲线。例如，设想在 $B=\R^2$ 中给定如图 4.2 所示交于一点 $b$ 的射线。取 $L_1,L_2,L_3$ 的指向点 $b$（即三条线段的交点）外侧的本原整切向量 $v_1,v_2,v_3\in\R^2$。现在我们有三个圆柱面 $X(L_i)$，它们在 $b$ 上方并不吻合：纤维 $f^{-1}(b)=\R^2/\Z^2$ 与 $X(L_i)$ 的交是一个圆 $\R v_i/\Z v_i$。这些圆在 $H_1(f^{-1}(b),\Z)=\Lambda_b$ 中恰好由向量 $v_1,v_2,v_3$ 代表，于是这些圆在 $f^{-1}(b)$ 中作为一个曲面的边界这一条件就是 $v_1+v_2+v_3=0$。因此，若这一条件成立，我们就可以粘入一个包含于 $f^{-1}(b)$ 中的曲面 $S$，使得 $X(L_1)\cup X(L_2)\cup X(L_3)\cup S$ 在 $b$ 处不再有边界。当然，它远非全纯子流形。然而预期是，这类对象可以被形变到附近的一条全纯曲线。
% [Tutor Note: 平衡条件 v_1+v_2+v_3=0 中的三个向量都落在 Λ_b 里（不是整体的 Λ），且方向约定是「从 b 指向外侧」，故和为零而不是 −v_1+⋯=0；文字层把 Rv_i/Zv_i 抽成 Rvi/Zvi，尚可辨认。]

确切地说，继续上面的例子，设 $b=0$、$v_1=(1,0)$、$v_2=(0,1)$ 以及 $v_3=(-1,-1)$。取 $X(B)=(\C^*)^2$ 上的全纯坐标 $q_1,q_2$，考虑由 $1+q_1+q_2=0$ 定义的曲线 $C\subseteq(\C^*)^2$。看这条曲线在映射 $f:X(B)\to B$ 下的像；这个映射在这里可以显式地写成 $(q_1,q_2)\mapsto-\frac{1}{2\pi}(\log|q_1|,\log|q_2|)$。可以发现，所得的是上述三价图的一个加厚，通常称为 amoeba（变形虫）。此外，若考虑的不是 $X(B)$ 而是 $X_\epsilon(B)$，此时全纯坐标由 $q_j=e^{2\pi i(x_j+iy_j)/\epsilon}$ 给出，而 $f_\epsilon:X_\epsilon(B)\to B$ 由 $(q_1,q_2)\mapsto-\frac{\epsilon}{2\pi}(\log|q_1|,\log|q_2|)$ 给出，则可发现当 $\epsilon\to0$ 时，$f_\epsilon(C)$ 收敛到上图中的三价图。在这个意义上，$B$ 上的三价图是一族趋于大复结构极限的流形上曲线的极限版本。

% [原文 p.16]
\figimg[0.58]{p4_fig02.png}
\figtitle{4.2}{}
% [Tutor Note: 原文图注仅有编号 "Figure 2."，无描述文字，故图题第二参数留空；图中三条射线 L_1,L_2,L_3 交于原点，对应 p15 的三价图与 v_1+v_2+v_3=0。]

代数环面中曲线的这一基本图景如今已被研究得非常透彻。特别地，这一研究催生了热带几何这门学科。「热带」一词的动机来自热带半环所起的作用。这个半环是 $(\R,\oplus,\odot)$，其中加法与乘法由
\begin{align*}
a\oplus b&:=\min(a,b),\\
a\odot b&:=a+b
\end{align*}
给出。
% [Tutor Note: 热带半环的「加法」取 min、「乘法」取 +，与通常算术正好「互换」；注意 ⊕ 与 ⊙ 都是自定义运算符，不是普通的圈加/圈乘，圆点符号是 \odot。]

「热带」一词用于纪念巴西数学家 Imre Simon，他开创了这个半环的使用。

我们现在如下考虑热带半环上的多项式。设 $S\subseteq\Z^n$ 是一个有限子集，考虑 $\R^n$ 上形如
\[
g:=\sum_{(p_1,\dots,p_n)\in S}c_{p_1\cdots p_n}x_1^{p_1}\cdots x_n^{p_n}
\]
的热带多项式，其中系数属于 $\R$，运算在热带半环中进行。则 $g$ 是 $\R^n$ 上的凸分段线性函数，而 $g$ 非线性的轨迹称为\emph{热带超曲面}。特别地，在 $n=2$ 的情形，我们得到一条热带曲线。在图 4.2 的例子中，相应的热带多项式可以取为 $0\oplus x_1\oplus x_2$。
% [Tutor Note: 热带多项式中「系数乘单项式」变成普通加法，故 0⊕x_1⊕x_2 即 min(0,x_1,x_2)，其非线性轨迹正是图 4.2 的三条射线。文字层把下标 c_{p_1…p_n} 抽成 cp1...pn，须按页图。]

尽管热带半环已在热带几何中被广泛使用，但对我们而言，用方程来理解 $B$ 上的热带曲线并不太方便，因为通常这些曲线会具有很高的余维数。因此，更好的做法是遵循 Mikhalkin [62]，使用参数化的热带曲线。

参数化热带曲线的定义域将是一个带权图。下文中，$\overline{\Gamma}$ 将表示一个连通图。这样的图可以从两种不同的方式来看。首先，它可以被看作一个纯组合对象，即一个顶点集 $\overline{\Gamma}^{[0]}$ 与一个边集 $\overline{\Gamma}^{[1]}$，后者由 $\overline{\Gamma}^{[0]}$ 中元素的无序对组成，表示一条边的两个端点。
% [原文 p.17]
我们也可以把 $\overline{\Gamma}$ 看作该图的拓扑实现，即一个拓扑空间，它是与各边对应的线段之并。我们将随意混用这两种观点。此后我们用 $\Gamma$ 记从 $\overline{\Gamma}$ 删去 $\overline{\Gamma}$ 的单价顶点后所得的拓扑空间，于是 $\Gamma$ 可能有一些非紧的边。

我们还要求 $\overline{\Gamma}$ 带有一个\emph{权函数}，即一个映射
\[
w:\overline{\Gamma}^{[1]}\to\N=\{0,1,2,\dots\}.
\]

把 $\R^n$ 换成一个一般的热带仿射流形 $B$，我们现在得到如下定义：

\begin{definition}{4.4.1}
$B$ 中一条\emph{参数化热带曲线}是一个连续映射
\[
h:\Gamma\to B,
\]
其中 $\Gamma$ 由如上所述的图 $\overline{\Gamma}$ 得到，并满足以下两条性质：
\begin{enumerate}
\item 若 $E\in\Gamma^{[1]}$ 且 $w(E)=0$，则 $h|_E$ 为常数；否则 $h|_E$ 是把 $E$ 作为具有有理斜率的线段、射线或直线映入 $B$ 的一个真嵌入。
\item \emph{平衡条件。}设 $V\in\overline{\Gamma}^{[0]}$ 是一个价数大于 $1$ 的顶点，相邻的边为 $E_1,\dots,E_\ell$。设 $v_i\in\Lambda_{h(V)}$ 是 $h(E_i)$ 在 $h(V)$ 处的本原切向量，方向背离 $h(V)$。则
\[
\sum_{i=1}^{\ell}w(E_i)v_i=0.
\]
\end{enumerate}
\end{definition}
% [Tutor Note: 定义中 h 的定义域是删去单价顶点后的 Γ，而平衡条件里的顶点 V 取在 Γ̄^{[0]} 中；两个符号 Γ 与 Γ̄ 不可混用（Γ̄ 才是带权函数 w 的原始图）。]

这里平衡条件只是在表达如下拓扑要求：各圆柱面 $X(h(E_i))\subseteq X(B)$ 的边界可以在 $X(B)\to B$ 在 $h(V)$ 上的纤维中用一个曲面连接起来。权可以解释为把圆柱面 $X(h(E_i))$ 以重数计入。

于是出现一个重要问题：

\begin{question}{4.4.2}
给定的参数化热带曲线何时可以被看作 $\epsilon\to0$ 时 $X_\epsilon(B)$ 中全纯曲线的极限？
\end{question}

这个问题在 $B=\R^n$ 时引起了大量关注，在 $n=2$ 的情形有完全令人满意的结果（答案：总是可以），而在 $n\ge3$ 时结果不那么完整。$n=2$ 的情形最早由 Mikhalkin [62] 处理，而所有维数的结果最早由 Nishinou 与 Siebert [64] 得到。特别地，Mikhalkin 证明了在这个二维情形，人们可以计算给定次数与亏格、经过一组固定点的曲线数目，这表明困难的全纯枚举问题可以纯组合地解决。这项工作给人希望：在远为一般的设定中真的可以用组合方式数曲线。再次地，在二维情形，我自己的工作 [24] 表明镜像一侧（对 $\PP^2$ 的镜像对称）
% [原文 p.18]
\figimg[0.62]{p4_fig03.png}
\figtitle{4.3}{}
% [Tutor Note: 原文图注仅有编号 "Figure 3."，无描述文字，故图题第二参数留空。图中画出二维底 B 中带孤立奇点处、沿从奇点出发的线段纤维化的全纯圆盘（上：三个带薄片的环面纤维；下：底 B 上的判别轨迹点与 X 标记）。]
也可以被热带地解释，从而给出 $\PP^2$ 的镜像对称的一个完全热带的解释。

至此我们还没有考虑 $B$ 带奇点的情形。当 $B$ 有奇点时，我们预期当某个顶点落入奇异轨迹的某一点内部时应当可以放松平衡条件；特别地，可以允许单价顶点映射到奇异轨迹。理由如下：一旦我们把 $X(B_0)$ 紧化为 $X(B)$，就预期会看到沿从奇点出发的线段纤维化的全纯圆盘：当 $B$ 是二维且具有孤立奇点时，图 4.3 描绘了这一情形。
% [Tutor Note: 「放松平衡条件」的具体含义是：允许单价顶点映射到奇异轨迹（判别轨迹）之内，此类顶点处不再要求 Σ w(E_i)v_i = 0；这并非删去平衡条件，而只是在奇点处允许例外，正文照原文措辞保留。]

我们将避免给出 $B$ 带奇点时热带曲线应当意味着什么的精确定义，主要因为目前还不清楚精确定义应该是什么。不过，希望上面的讨论已经足够清楚：在直观层面，数曲线应当是可以在 $B$ 上完成的事情。

\section{SYZ 猜想的问题，以及如何绕开它们}
第 3 节的讨论表明，SYZ 猜想在纯拓扑层面给出了镜像对称的一个漂亮描述。这一点本身往往就很有用，但它无法触及原本困难的微分几何猜想，也无法说明为什么镜像对称能数曲线。
